\section{What is not rigorous here}

\begin{enumerate}
\item FNR-gate convergence (heuristic). We assume each round returns a
family meeting the $(\mathrm{fpr},\mathrm{fnr})$ budget and that the boundary
feedback strictly reduces the unresolved set. The clustering
(\texttt{identify\_cluster\_around}) is a greedy heuristic with no proven
guarantee that it finds a coherent family when one exists; the \texttt{fnr\_limit}
check certifies the result but not that the search reaches it. The stall
detector exists precisely because this can fail to make progress.

\item Sequential peeking in the split test. Lemma~\ref{lem:split-sound}
is stated for a fixed member sample. Reading members until the Bayes factor
crosses the threshold (an optional-stopping rule) inflates the false-split rate
above \texttt{split\_pval} unless the threshold is replaced by an anytime-valid
one. The current code reads a fixed cap of members, so it is sound (the fixed-$n$
threshold is valid) rather than sequential.

\item Coverage of transient states. Assumption~\ref{ass:cover} is a real
restriction: fixed-length sampling can miss states reachable only off the sampled
length, and the guarantee is only over $D$, not worst-case.

\item Independence approximations. Rows share cells (a prefix's noise
bits recur across families across rounds), so the ``independent noise'' used in
the union bounds is an approximation once the same strings are reused.
\end{enumerate}
